\section{RQ1: Integration Architectures}
\label{sec:taxonomy}

% Schematic of the seven integration architectures (A1--A7).
\begin{figure*}[t]
\centering
\begin{tikzpicture}[
  font=\scriptsize,
  lab/.style={align=left, anchor=west, text width=33mm},
  lm/.style={draw=nsblue, fill=nsblue!10, rounded corners=2pt, minimum height=6mm, minimum width=13mm, align=center, line width=0.6pt},
  sym/.style={draw=A3c!80!black, fill=A3c!12, rounded corners=2pt, minimum height=6mm, minimum width=15mm, align=center, line width=0.6pt},
  dat/.style={align=center, text=black!75, font=\tiny, inner sep=1pt},
  arr/.style={-{Stealth[length=1.4mm]}, line width=0.5pt},
  fb/.style={-{Stealth[length=1.4mm]}, line width=0.5pt, dashed, draw=A5c},
  cnt/.style={anchor=east, font=\scriptsize\bfseries, text=nsblue}]
\def\rowgap{11.5mm}
\def\xs{41mm}
% ---- A1
\node[lab] (l1) at (0,0) {\textbf{A1}\enspace Translate-and-solve};
\node[dat] (a1i) at (\xs,0) {problem\\(text)};
\node[lm, right=5mm of a1i] (a1m) {LM};
\node[dat, right=5mm of a1m] (a1f) {formal\\spec};
\node[sym, right=5mm of a1f] (a1s) {solver /\\prover};
\node[dat, right=5mm of a1s] (a1o) {answer};
\draw[arr] (a1i)--(a1m); \draw[arr] (a1m)--(a1f); \draw[arr] (a1f)--(a1s); \draw[arr] (a1s)--(a1o);
\node[cnt] at (16.3,0) {$n=\cnt{A1}$};
% ---- A2
\node[lab] (l2) at (0,-\rowgap) {\textbf{A2}\enspace Verifier-in-the-loop};
\node[dat] (a2i) at (\xs,-\rowgap) {task};
\node[lm, right=5mm of a2i] (a2m) {LM};
\node[dat, right=5mm of a2m] (a2f) {candidate\\proof/code};
\node[sym, right=5mm of a2f] (a2s) {checker};
\node[dat, right=5mm of a2s] (a2o) {verified\\output};
\draw[arr] (a2i)--(a2m); \draw[arr] (a2m)--(a2f); \draw[arr] (a2f)--(a2s); \draw[arr] (a2s)--(a2o);
\draw[fb] (a2s.north) -- ++(0,2.2mm) -| node[pos=0.25, above, font=\tiny, text=A5c, inner sep=1pt] {verdict / error feedback} (a2m.north);
\node[cnt] at (16.3,-\rowgap) {$n=\cnt{A2}$};
% ---- A3
\node[lab] (l3) at (0,-2*\rowgap) {\textbf{A3}\enspace Symbolic guidance\\\phantom{\textbf{A3}}\enspace or constraints};
\node[dat] (a3i) at (\xs,-2*\rowgap) {task};
\node[lm, right=5mm of a3i] (a3m) {LM};
\node[dat, right=5mm of a3m] (a3o) {output};
\node[sym, right=8mm of a3o] (a3s) {grammar, search,\\constraints, loss};
\draw[arr] (a3i)--(a3m); \draw[arr] (a3m)--(a3o);
\draw[arr, draw=A3c!80!black] (a3s.south) -- ++(0,-2.2mm) -| node[pos=0.25, below, font=\tiny, text=A3c!70!black, inner sep=1pt] {steers search, decoding or training} (a3m.south);
\node[cnt] at (16.3,-2*\rowgap) {$n=\cnt{A3}$};
% ---- A4
\node[lab] (l4) at (0,-3*\rowgap) {\textbf{A4}\enspace LM-induced\\\phantom{\textbf{A4}}\enspace symbolic artifacts};
\node[dat] (a4i) at (\xs,-3*\rowgap) {traces,\\data};
\node[lm, right=5mm of a4i] (a4m) {LM};
\node[dat, right=5mm of a4m] (a4f) {rules, programs,\\predicates};
\node[sym, right=5mm of a4f] (a4s) {library /\\executor};
\node[dat, right=5mm of a4s] (a4o) {reused across\\tasks};
\draw[arr] (a4i)--(a4m); \draw[arr] (a4m)--(a4f); \draw[arr] (a4f)--(a4s); \draw[arr] (a4s)--(a4o);
\node[cnt] at (16.3,-3*\rowgap) {$n=\cnt{A4}$};
% ---- A5
\node[lab] (l5) at (0,-4*\rowgap) {\textbf{A5}\enspace Symbolic-to-neural\\\phantom{\textbf{A5}}\enspace transfer};
\node[sym] (a5s) at ($(\xs,-4*\rowgap)+(6mm,0)$) {symbolic\\procedure};
\node[dat, right=5mm of a5s] (a5f) {training data,\\targets};
\node[lm, right=5mm of a5f] (a5m) {LM};
\node[dat, right=5mm of a5m] (a5o) {trained /\\distilled model};
\draw[arr] (a5s)--(a5f); \draw[arr] (a5f)--(a5m); \draw[arr] (a5m)--(a5o);
\node[cnt] at (16.3,-4*\rowgap) {$n=\cnt{A5}$};
% ---- A6
\node[lab] (l6) at (0,-5*\rowgap) {\textbf{A6}\enspace Evaluation with a\\\phantom{\textbf{A6}}\enspace symbolic oracle};
\node[sym] (a6s) at ($(\xs,-5*\rowgap)+(6mm,0)$) {solver /\\generator};
\node[dat, right=5mm of a6s] (a6f) {problems +\\gold answers};
\node[lm, right=5mm of a6f] (a6m) {LM under\\test};
\node[dat, right=5mm of a6m] (a6o) {scores,\\failure modes};
\draw[arr] (a6s)--(a6f); \draw[arr] (a6f)--(a6m); \draw[arr] (a6m)--(a6o);
\node[cnt] at (16.3,-5*\rowgap) {$n=\cnt{A6}$};
% ---- A7
\node[lab] (l7) at (0,-6*\rowgap) {\textbf{A7}\enspace LM as interface to\\\phantom{\textbf{A7}}\enspace a symbolic core};
\node[dat] (a7i) at (\xs,-6*\rowgap) {image, video,\\text};
\node[lm, right=5mm of a7i] (a7m) {LM / VLM};
\node[dat, right=5mm of a7m] (a7f) {grounded\\symbols};
\node[sym, right=5mm of a7f] (a7s) {symbolic\\core};
\node[dat, right=5mm of a7s] (a7o) {decision,\\explanation};
\draw[arr] (a7i)--(a7m); \draw[arr] (a7m)--(a7f); \draw[arr] (a7f)--(a7s); \draw[arr] (a7s)--(a7o);
\node[cnt] at (16.3,-6*\rowgap) {$n=\cnt{A7}$};
% legend
\node[lm, minimum width=8mm, minimum height=4mm, font=\tiny] (k1) at (\xs,-6.8*\rowgap) {LM};
\node[right=1mm of k1, font=\tiny] (k1t) {neural (language or vision-language model)};
\node[sym, minimum width=8mm, minimum height=4mm, font=\tiny, right=4mm of k1t] (k2) {sym};
\node[right=1mm of k2, font=\tiny] {symbolic component};
\end{tikzpicture}
\caption{The seven integration architectures identified in the \cnt{included} included studies, with the number of studies coded under each. Blue boxes are learned models; green boxes are symbolic components. Dashed arrows denote feedback from the symbolic component to the model.}
\label{fig:architectures}
\end{figure*}


The \cnt{included} studies fall into seven architectures (Fig.~\ref{fig:architectures}). Checker-in-the-loop designs dominate (A2, \cnt{A2} studies, \cnt{A2pct}\%), followed by symbolic-oracle evaluations (A6, \cnt{A6}), translate-and-solve pipelines (A1, \cnt{A1}) and symbolic guidance (A3, \cnt{A3}). Symbolic-to-neural transfer (A5, \cnt{A5}), model-induced artifacts (A4, \cnt{A4}) and model-as-interface designs (A7, \cnt{A7}) are less common. Across all studies, the symbolic side gives an exact verdict on the final output in \cnt{exact} studies, an empirical check in \cnt{empirical}, and no check in \cnt{none}. The architectures differ in where trust is placed, which determines what a system can guarantee.

\textbf{A2: Checker-in-the-loop (\cnt{A2}).} A checker or executor returns verdicts that select, repair, filter or reward model outputs. The category splits by check strength.
\begin{itemize}
  \item \emph{Exact} (\cnt{A2exact} studies). Here the proof assistant's acceptance is a guarantee relative to the formal statement. Typical approaches are proof search~\cite{lample2022hypertree,jiang2022thor,yang2023leandojo,xin2025bfs}, proof-assistant feedback for reinforcement learning~\cite{xin2025deepseek,dong2025stp,lin2026goedel}, and draft-and-prove pipelines~\cite{jiang2023draft,cao2025reviving,varambally2026hilbert}. Verification-aware code generation in Dafny or Verus~\cite{aggarwal2025alphaverus,yang2026exverus}, model-checking repair loops~\cite{agrawal2025specmas} and SAT-certified reasoning rewards~\cite{liu2025saturn} belong here too.
  \item \emph{Empirical} (\cnt{A2empirical} studies). Here the check can pass a wrong output: execution feedback for code and programming by example~\cite{ni2023lever,wang2024hypothesis,lee2025program}, solver feedback in optimisation modelling~\cite{ahmaditeshnizi2024optimus,chen2025solver}, data fit in LLM-guided equation search~\cite{shojaee2025sr,grayeli2024symbolic}, and simulator or environment success in embodied agents~\cite{ma2024eureka,mahdavi2024leveraging}.
\end{itemize}
Because A2 thus mixes sound verification with tests and data fit, we report the two kinds separately and restrict statements about guarantees to exact checks.

\textbf{A6: Symbolic-oracle evaluation (\cnt{A6}).} Benchmarks whose problems are generated, or whose answers are certified, by a symbolic system. Examples are solver-generated logic puzzles~\cite{lin2025zebralogic,wei2025satbench}, prover-checked first-order reasoning~\cite{qi2025large,han2024folio}, and planning benchmarks with formal validators~\cite{zuo2025planetarium,kokel2025acpbench}. Formal-mathematics and verification benchmarks~\cite{zheng2022minif2f,tsoukalas2024putnambench,hu2025minictx,ye2026verina} and symbolic perturbations of mathematical word problems~\cite{mirzadeh2025gsm} also belong here. Most A6 studies certify answers empirically (\cnt{A6empirical}), for example by answer matching against a generator. Only \cnt{A6exact} use a sound checker.

\textbf{A1: Translate-and-solve (\cnt{A1}).} The model writes a formal representation once and an engine computes the answer, so correctness reduces to the faithfulness of the translation. Program-aided reasoning~\cite{gao2023pal,gou2024tora,wang2024mathcoder} and binding language models to symbolic languages~\cite{cheng2023binding} are the largest group. Others translate text to first-order logic or SAT~\cite{olausson2023linc,ye2023satlm,kesseli2025logic}, answer-set programs~\cite{alviano2025integrating,santana2024automatic}, planning languages~\cite{hao2025planning,huang2025planning} or optimisation models~\cite{astorga2025autoformulation,yang2025optibench}. By construction, A1 pipelines do not check the model's formalisation (all \cnt{A1none} are coded \texttt{none}), which is why reporting formalisation faithfulness matters (Section~\ref{sec:evaluation}).

\textbf{A3: Symbolic guidance (\cnt{A3}).} Symbolic structure steers generation without judging it. Examples include constrained decoding with grammars and automata~\cite{scholak2021picard,wang2023grammar,beurerkellner2024guiding,banerjee2025crane}, logical decoding constraints~\cite{lu2021neurologic,lu2022neurologic,zhang2023tractable}, and static or dataflow analysis that guides code models~\cite{mukherjee2021neural,wei2023typet5,li2025dataflow}. Context-free surrogates of model proposals that guide enumerative synthesis~\cite{barke2024hysynth} and neurosymbolic training losses~\cite{ahmed2023pseudo} also belong here.

\textbf{A5: Symbolic-to-neural transfer (\cnt{A5}).} Symbolic procedures generate training data or signal offline, and the model runs without them: synthetic logic corpora~\cite{clark2020transformers,morishita2023learning,morishita2024enhancing}, formal theorem--proof data~\cite{huang2024mustard,ying2024lean,wu2025alchemy}, solver-generated mathematics~\cite{li2024neuro}, interpreter-generated execution traces~\cite{li2025codeio,wang2026stepcodereasoner}, and consistency losses against a rule base~\cite{calanzone2025logically}.

\textbf{A4: LM-induced artifacts (\cnt{A4}).} The model writes rules, programs, predicates or reward machines that are stored and reused, but no symbolic checker filters them; libraries that are filtered by tests or verification are A2 under the decision tree. Examples include programmatic world models~\cite{tang2024worldcoder,piriyakulkij2025poe}, invented predicates and planning domains for robots~\cite{liang2025visualpredicator,liang2026exopredicator,ye2025unidomain}, interpretable rule libraries~\cite{wang2025large,zhang2025ruag}, and reward functions or reward machines~\cite{xie2024text2reward,castanyer2026arm}.

\textbf{A7: LM as interface (\cnt{A7}).} A language or vision-language model grounds perception or produces explanations, and a symbolic core decides. Examples include logic-rule latent models over vision-language judgements~\cite{he2026unveiling}, robustness of vision-language reasoning with a symbolic executor~\cite{chen2026vlms}, programmable neurosymbolic learning frameworks~\cite{naik2025dolphin}, and knowledge-enhanced logical guardrails~\cite{kang2025r2}.

\subsection{Relation to existing integration taxonomies}
\label{sec:taxonomy-relations}
Table~\ref{tab:taxmap} relates the architectures to Kautz's integration types~\cite{kautz2022third} and to the perspectives of Yang et al.~\cite{yang2025neuro}. The correspondence is approximate, but it shows what the finer scheme adds. Kautz's Neuro|Symbolic pipeline covers A1, A4 and A7, which differ in what the symbolic side receives: a whole problem, a reusable artifact or grounded percepts. Yang et al.'s LLM+Symbolic perspective merges A2, where the symbolic side judges, with A3, where it only steers. The check-strength attribute adds a second distinction no earlier scheme records: whether that judgement is sound.

\begin{table}[t]
\centering
\caption{Approximate correspondence with existing taxonomies.}
\label{tab:taxmap}
\footnotesize
\begin{tabularx}{\columnwidth}{@{}l X X@{}}
\toprule
Ours & Kautz~\cite{kautz2022third} & Yang et al.~\cite{yang2025neuro} \\
\midrule
A1 & Neuro|Symbolic & LLM$\rightarrow$Symbolic \\
A2 & Neuro|Symbolic with feedback & LLM+Symbolic \\
A3 & Symbolic[Neuro] (search); Neuro$_\text{Symbolic}$ (loss) & LLM+Symbolic; Symbolic$\rightarrow$LLM \\
A4 & Neuro|Symbolic & LLM$\rightarrow$Symbolic \\
A5 & Neuro:Symbolic$\rightarrow$Neuro & Symbolic$\rightarrow$LLM \\
A6 & -- (evaluation) & -- \\
A7 & Neuro|Symbolic & LLM$\rightarrow$Symbolic \\
\bottomrule
\end{tabularx}
\end{table}
